% Adapted from DriftNeRF's images/fig-overview.tex: ray/sample geometry,
% contrasting anchor/candidate colors, proximity links, rounded consensus box.
% https://github.com/qqgjyx/overleaf_DriftNeRF/blob/79946b56b0b951b03da3b58f4ea7062e4076bfa2/images/fig-overview.tex
% Citation: Wang et al., arXiv:2610.29106v1 (CC BY 4.0).
% Changes: compact two-panel layout, input-ray labels, equal candidate counts,
% stratified off-ray queries, explicit detached targets. No result data.
% Blank lines are commented because acmart collects teaser content as an argument.
\begin{teaserfigure}
\centering
\begin{tikzpicture}[
  x=1cm,y=1cm,font=\sffamily\small,
  candidate/.style={circle,fill=blue!65!black,inner sep=2pt},
  anchorpoint/.style={circle,fill=orange!85!black,inner sep=2.6pt},
  arr/.style={-{Stealth[length=4pt]},semithick},
  target/.style={draw=black!45,rounded corners=2pt,fill=black!3,
    minimum height=0.55cm,align=center,inner sep=4pt}
]
\begin{scope}
  \node[font=\sffamily\bfseries] at (3.45,2.8)
    {(a) DriftNeRF: reuse ray samples};
  \draw[gray!40,fill=gray!3] (0.2,0) rectangle (6.7,2.25);
  \draw[orange!70!black,thin] (0.35,0.25) -- (6.55,2.0)
    node[pos=0.05,above,font=\scriptsize] {ray $a$};
  \draw[blue!55!black,thin] (0.35,1.75) -- (6.55,1.75)
    node[pos=0.97,above,font=\scriptsize] {ray $b$};
  \draw[blue!55!black,thin] (0.35,0.5) -- (6.55,0.5)
    node[pos=0.97,above,font=\scriptsize] {ray $c$};
  \coordinate (a) at (3.45,1.125);
  \foreach \x/\y in {1.5/1.75,4.3/1.75,2.4/0.5,5.4/0.5}{
    \draw[black!35,densely dotted] (a) -- (\x,\y);
    \node[candidate] at (\x,\y) {};
  }
  \node[anchorpoint,label={[font=\scriptsize]below:$\vx_i$}] at (a) {};
  \node[font=\scriptsize,align=center] at (3.45,-0.38)
    {Other-ray candidates; original distance cutoff};
  \node[target] (leftTarget) at (3.45,-1.17)
    {Distance-weighted average of $\sg(\sigma(\vx_j))$};
  \draw[arr,black!55] (3.45,-0.62) -- (leftTarget.north);
\end{scope}
\begin{scope}[xshift=8.1cm]
  \node[font=\sffamily\bfseries] at (3.45,2.8)
    {(b) Proposed: query local cells};
  \draw[gray!40,fill=gray!3] (0.2,0) rectangle (6.7,2.25);
  \draw[gray!35] (3.45,0) -- (3.45,2.25);
  \draw[gray!35] (0.2,1.125) -- (6.7,1.125);
  \coordinate (a) at (3.45,1.125);
  \foreach \x/\y in {1.3/1.65,4.6/1.85,2.2/0.4,5.65/0.65}{
    \draw[black!35,densely dotted] (a) -- (\x,\y);
    \node[candidate] at (\x,\y) {};
  }
  \node[anchorpoint,label={[font=\scriptsize]above:$\vx_i$}] at (a) {};
  \node[font=\scriptsize,align=center] at (3.45,-0.38)
    {One query per cell; no originating ray ID};
  \node[target] (rightTarget) at (3.45,-1.17)
    {Distance-weighted average of $\sg(\sigma(\vx_j))$};
  \draw[arr,black!55] (3.45,-0.62) -- (rightTarget.north);
\end{scope}
\end{tikzpicture}
\caption{Sampling constructions, adapted from DriftNeRF~\citep{wang2026driftnerf}. Orange denotes the anchor; blue denotes four candidates per panel. Dotted links connect candidates to the anchor for distance-weighted averaging. This conceptual 2D slice illustrates 3D sampling; equal displayed counts do not imply equal query costs.}
\label{fig:sampling-hypothesis}
\Description{Two panels with equal candidate counts. Left: an anchor on input ray a pools four candidate densities from other input rays b and c. Right: an anchor queries one candidate in each of four local cells, with no ray identifiers. Dotted proximity links lead to a distance-weighted detached-density target in each panel.}
\end{teaserfigure}
